\chapter{Strings, formats and JSON}
\label{ch:strings}

An agent's life is text: prompts, tool output, file contents, the model's
JSON. This chapter covers the string toolkit, the two ways to print, a
string literal made for prompts, and how Scheme and JSON map onto each
other.

\section{The string toolkit}

\begin{scm}
(string-append "tom" "ato")                  ;=> "tomato"
(string-length "ñandú")                      ;=> 5      characters, not bytes
(substring "watering can" 0 8)               ;=> "watering"
(string-upcase "basil")                      ;=> "BASIL"
(string-trim "  mint \n")                    ;=> "mint"
(string-split "tomato,basil,mint" ",")       ;=> ("tomato" "basil" "mint")
(string-join '("a" "b" "c") ", ")            ;=> "a, b, c"
(string-contains? "raised bed" "bed")        ;=> #t
(string-prefix? "go " "go test ./...")       ;=> #t
(string-suffix? ".scm" "init.scm")           ;=> #t
(string-replace "slug slug" "slug" "snail")  ;=> "snail snail"
(string<? "apple" "banana")                  ;=> #t
(string->symbol "basil")                     ;=> basil
(string->number "3.5")                       ;=> 3.5
(string->number "12abc")                     ;=> #f
\end{scm}

Two things to notice. Indices and lengths count \emph{characters}, so
accented text behaves. And \code{string-split} splits on a literal
separator, not a pattern: there are no regular expressions in the
language. When you need one, the \code{index-grep} primitive runs a real
one over files (Chapter~\ref{ch:primitives}).

\begin{careful}
Look at the argument order of \code{string-prefix?} and
\code{string-suffix?}: the \emph{prefix comes first}, then the string. It
reads as ``is \code{"go "} a prefix of \code{"go test ./..."}?''.
\code{string-contains?} is the other way round: the haystack, then the
needle.
\end{careful}

\section{\texttt{display} and \texttt{write}}

There are two ways to turn a value into text. \code{display} is for
people: strings without their quotes. \code{write} is for machines:
strings quoted and escaped, so the text reads back as the same value.

\begin{scm}
(with-output-to-string
  (display "basil") (newline) (write "basil"))
                                ;=> "basil\n\"basil\""
\end{scm}

\code{with-output-to-string} runs its body and returns everything the body
printed, instead of printing it. It is the language's one string port, and
it is surprisingly handy for building reports piece by piece.

\section{\texttt{format}}

\code{format} fills a template. It knows four directives:

\begin{center}
\begin{tabular}{ll}
\toprule
\code{\textasciitilde a} & the next argument, as \code{display} shows it \\
\code{\textasciitilde s} & the next argument, as \code{write} shows it \\
\code{\textasciitilde\%} & a newline \\
\code{\textasciitilde\textasciitilde} & a tilde \\
\bottomrule
\end{tabular}
\end{center}

\begin{scm}
(format "~a has ~a leaves" "basil" 12)       ;=> "basil has 12 leaves"
(format "~s" "quoted")                       ;=> "\"quoted\""
(format "~a" '(1 "two" three))               ;=> "(1 two three)"
(format "~s" '(1 "two" three))               ;=> "(1 \"two\" three)"
(format "~a ~a" 1)       ; error: format: too few arguments for "~a ~a"
\end{scm}

Too few or too many arguments is an error, which catches a lot of bugs in
prompt templates. \code{(format "\textasciitilde a" x)} is also the quickest way to turn
\emph{anything} into a string, which \nodwyn's own Scheme does all the time
when a value might be a number, a string or a symbol.

\section{Raw strings: prompts without backslashes}

Prompts are long, full of quotes and backticks, and written by people. For
them, \nodwyn{} Scheme has a \textbf{raw string}: three double quotes, no
escapes at all.

\begin{scm}
(define pattern """\d+\.\d+""")
pattern                              ;=> "\\d+\\.\\d+"
"""a "b" c"""                        ;=> "a \"b\" c"
\end{scm}

When the opening \code{"""} is followed directly by a newline, the block is
read as \textbf{indented text}: the shared indentation is removed, along
with the first and last newlines. So you can indent a prompt with your code
and get text with no indentation at all:

\begin{scm}
(define instructions
  """
  Prefer table-driven tests.
    Indented further keeps the difference.
  Run `make accept` before saying a change is done.
  """)
\end{scm}

\code{instructions} is exactly:

\begin{txt}
Prefer table-driven tests.
  Indented further keeps the difference.
Run `make accept` before saying a change is done.
\end{txt}

You will use this in Chapter~\ref{ch:checks} to extend the system prompt.

\section{JSON}

Two builtins cross between JSON text and Scheme values:

\begin{scm}
(json-parse "{\"plant\": \"basil\", \"height\": 12.5,
              \"tags\": [\"herb\", \"sun\"], \"watered\": true, \"note\": null}")
;=> (("plant" . "basil") ("height" . 12.5) ("tags" "herb" "sun")
;    ("watered" . #t) ("note" . null))

(json-emit (list (cons "plant" "basil")
                 (cons "tags" (list "herb" "sun"))))
;=> "{\"plant\":\"basil\",\"tags\":[\"herb\",\"sun\"]}"
\end{scm}

\begin{center}
\begin{tabular}{ll}
\toprule
\textbf{JSON} & \textbf{Scheme} \\
\midrule
object & association list with \emph{string} keys, in order \\
array & list \\
string, number & string, integer or float \\
\code{true}, \code{false} & \code{\#t}, \code{\#f} \\
\code{null} & the symbol \code{null} \\
\bottomrule
\end{tabular}
\end{center}

Why does \code{("tags" "herb" "sun")} have no dot? Because it \emph{is}
\code{("tags" . ("herb" "sun"))}: a pair whose \code{cdr} is a list prints
as one longer list. \code{(assoc-ref parsed "tags")} still gives you
\code{("herb" "sun")}.

Going the other way, \code{json-emit} treats a list whose elements are all
pairs with string keys as an object, and any other list as an array. One
consequence: \code{'()} emits \code{[]}, and there is no way to write an
empty object \code{\{\}}.

This mapping matters because it is exactly how a model's tool-call
arguments reach Scheme. When the model calls
\code{fs-read} with \code{\{"path": "plants.txt"\}}, \nodwyn{} parses the
JSON with these rules and applies \code{fs-read} to the values in the order
the tool declares.

\begin{puzzle}[ 6: a label maker]
Write \code{(label plant days)} that returns \code{"BASIL: water every 2
days"}, and \code{"MINT: water every day"} when \code{days} is 1. Use
\code{format} and \code{string-upcase}.\ans{6}
\end{puzzle}

\begin{puzzle}[ 7: CSV to JSON]
Write \code{(csv->json text)} that turns

\code{"plant,water\textbackslash nbasil,daily\textbackslash ntomato,weekly"}

into the JSON text
\begin{txt}
[{"plant":"basil","water":"daily"},{"plant":"tomato","water":"weekly"}]
\end{txt}
The first line holds the keys.\ans{7}
\end{puzzle}
